\section{World Model RL Training}\label{appendix:world_model}

This appendix reports on training PPO entirely inside a learned world model,
replacing the physics simulator with a data-driven dynamics model.

\subsection{World Model Architecture}

The world model is a 2-layer LSTM~\cite{hochreiter1997lstm}
(hidden dim $= 64$, sequence length $= 10$), implemented in
PyTorch~\cite{paszke2019pytorch} and trained on a 1\,M-transition
random-exploration dataset ($\sim$14 hours of operation; collection protocol
in Appendix~\ref{appendix:dataset_details}). Its input is the
concatenation of the current state
$\mathbf{s}_t = [x_c, \dot{x}_c, \theta_b, \dot{\theta}_b]$ and the action
$a_t$, and it predicts the state delta
$\Delta\mathbf{s}_{t+1} = \mathbf{s}_{t+1} - \mathbf{s}_t$, with normalization
statistics precomputed from the dataset (Fig.~\ref{fig:wm_arch}). We train it
for 50 epochs, reaching a final training-set prediction MAE of 0.611 in the
normalized state space.

\paragraph{Training hyperparameters.}
The model (hidden dim 64, 2 LSTM layers, dropout 0.2) is trained with the
Adam optimizer at learning rate $10^{-4}$, batch size 256, over
10-step sequences for 50 epochs. The objective is a per-dimension-weighted MSE
on the predicted state delta $\Delta\mathbf{s}$, with weights $[1, 5, 1, 5]$
on $[x_c, \dot{x}_c, \theta_b, \dot{\theta}_b]$ to emphasize the two velocity
channels. The 1\,M-transition collection is split $80/10/10$ into
train/validation/test, with normalization statistics computed on the training
split only, and all seeding uses seed~42. When MPPI plans with
this world model at deployment, we run it on CPU, which is faster than GPU for
a network this small on our target hardware, and the per-step planning cost
(about 40\,ms for 800 sampled rollouts over a horizon of 30) stays within the
50\,ms control period. PPO-WM instead deploys only its trained policy network,
so the world model runs online only under MPPI.

\begin{figure}[!t]
\centering
\resizebox{0.98\columnwidth}{!}{%
\begin{tikzpicture}[
    >=Stealth,
    every node/.style={font=\sffamily\small},
    box/.style={
        draw=black, thick, rounded corners, fill=white,
        minimum height=1.0cm, align=center, inner sep=5pt
    },
    flow/.style={draw=black, thick, -{Stealth}},
    skip/.style={draw=black, thick, dashed, -{Stealth}},
    node distance=0.9cm and 1.1cm
]
\definecolor{cInput}{RGB}{222,235,247}
\definecolor{cInputLine}{RGB}{91,143,185}
\definecolor{cLSTM}{RGB}{226,239,218}
\definecolor{cLSTMLine}{RGB}{112,160,90}
\definecolor{cHead}{RGB}{252,235,212}
\definecolor{cHeadLine}{RGB}{196,143,75}
\definecolor{cDelta}{RGB}{237,224,237}
\definecolor{cDeltaLine}{RGB}{150,110,150}
\definecolor{cOutput}{RGB}{222,235,247}
\definecolor{cOutputLine}{RGB}{91,143,185}

% Layer 1
\begin{scope}
  \node[box, fill=cLSTM!40, draw=cLSTMLine, minimum width=2.4cm,
        yshift=0.44cm] (lstm1b3) {};
  \node[box, fill=cLSTM!40, draw=cLSTMLine, minimum width=2.4cm,
        yshift=0.30cm] (lstm1b2) {};
  \node[box, fill=cLSTM!60, draw=cLSTMLine, minimum width=2.4cm,
        yshift=0.15cm] (lstm1b1) {};
  \node[box, fill=cLSTM, draw=cLSTMLine, thick, minimum width=2.4cm] (lstm1)
      {LSTM layer 1\\[-1pt]\tiny dim = 64};
\end{scope}

% Layer 2
\begin{scope}[xshift=3.2cm]
  \node[box, fill=cLSTM!40, draw=cLSTMLine, minimum width=2.4cm,
        yshift=0.44cm] (lstm2b3) {};
  \node[box, fill=cLSTM!40, draw=cLSTMLine, minimum width=2.4cm,
        yshift=0.30cm] (lstm2b2) {};
  \node[box, fill=cLSTM!60, draw=cLSTMLine, minimum width=2.4cm,
        yshift=0.15cm] (lstm2b1) {};
  \node[box, fill=cLSTM, draw=cLSTMLine, thick, minimum width=2.4cm] (lstm2)
      {LSTM layer 2\\[-1pt]\tiny dim = 64};
\end{scope}

\begin{scope}[on background layer]
\node[draw=cLSTMLine, thick, rounded corners, fit=(lstm1b2)(lstm2b2),
      inner sep=16pt, fill=cLSTM!15] (lstmcluster) {};
\end{scope}

% Inputs
\node[draw, rounded corners=2pt, inner sep=3pt, font=\sffamily\footnotesize,
      above=1.3cm of lstm1] (concat) {concat};
\node[box, fill=cInput, draw=cInputLine, minimum width=2.1cm,
      above=0.8cm of concat, xshift=-1.7cm] (action)
    {$a_t$\\[-1pt]\tiny action};
\node[box, fill=cInput, draw=cInputLine, minimum width=3.6cm,
      above=0.8cm of concat, xshift=1.8cm] (state)
    {$s_t = (x_c,\dot{x}_c,\theta_b,\dot{\theta}_b)$\\[-1pt]
     \tiny state};

% Output head
\node[box, fill=cHead, draw=cHeadLine, minimum width=2.1cm,
      right=1.5cm of lstm2] (head)
    {Linear\\head\\[-1pt]\tiny $64 \to 4$};
\node[box, fill=cDelta, draw=cDeltaLine, minimum width=2.5cm,
      right=1.1cm of head] (delta)
    {$\Delta s_{t+1}$\\[-1pt]\tiny state delta};

\node[fill=none, draw=none, font=\sffamily\Large,
      above=1.3cm of delta] (add2) {$\oplus$};
\node[box, fill=cOutput, draw=cOutputLine, minimum width=2.8cm,
      right=1.1cm of add2] (output)
    {$s_{t+1}$\\[-1pt]\tiny predicted next state};

% Edges
\draw[flow] (action) -- (concat);
\draw[flow] (state) -- (concat);
\draw[flow] (concat) -- (lstm1b3);
\draw[flow] (lstm1) -- (lstm2);
\draw[flow] (lstm2) -- (head);
\draw[flow] (head) -- (delta);
\draw[flow] (delta) -- (add2);
\draw[flow] (add2) -- (output);

\draw[skip] (state.east) -- ++(right:0.5cm)
    |- node[pos=0.85, above, font=\sffamily\tiny] {$s_t$} (add2.west);

\end{tikzpicture}%
}
\caption{World model architecture (v1).  Inputs $s_t$ (4-D state, z-score
  normalized) and $a_t$ (scalar action, used in its original bounded scale)
  are concatenated and fed through
  a 2-layer LSTM (hidden dim $= 64$, dropout $= 0.2$).  A linear output
  head maps to the 4-D state delta $\Delta s_{t+1}$; a skip connection
  adds this to the current state to produce the predicted next state
  $s_{t+1}$.}
\label{fig:wm_arch}
\end{figure}

\subsection{Checkpoint Provenance and Sign Conventions}
\label{appendix:wm_provenance}

MPPI and PPO-WM use \emph{separate} LSTM checkpoints of the same
architecture, trained on two distinct 1\,M-transition hardware
collections whose difference is subtle but consequential for RL
training. Both are released in the repository.

\begin{table}[H]
\centering
\caption{World model checkpoint provenance.}
\label{tab:wm_provenance}
\scriptsize
\setlength{\tabcolsep}{2pt}
\begin{tabular}{lp{3cm}p{3cm}}
\toprule
\textbf{Property} & \textbf{MPPI ckpt} & \textbf{PPO-WM ckpt} \\
\midrule
Dataset & \texttt{arcball\_\allowbreak cont\_\allowbreak 1M\_\allowbreak calib196\_\allowbreak 160.h5} & \texttt{arcball\_\allowbreak cont\_\allowbreak 1M\_\allowbreak pre\_\allowbreak neg.h5} \\
Sampling rate & 20\,Hz (50\,ms) & 20\,Hz (50\,ms) \\
Cart-velocity sign & negated (positive = right) & raw encoder (positive = left) \\
Action$\to$velocity $r$ & $+0.14$ (strong) & $\approx 0$ (negligible) \\
Sim SR of PPO-WM & $\sim$60\% (wall-stuck) & $\sim$98\% (robust) \\
Sim SR of MPPI & 100\% & N/A \\
\bottomrule
\end{tabular}
\end{table}

The two collections differ in their velocity-sign convention. The
MPPI dataset was collected with the current serial reader, which negates
the encoder velocity so that positive cart velocity = right in the
Python state space. The PPO-WM dataset was collected earlier on a
separate host whose reader stored the raw encoder velocity (positive =
left). The resulting model shows only a delayed, weak action response on the
cart-velocity channel: the per-step velocity change is dominated by encoder
noise, so the immediate effect of the action on cart velocity is small (a
response emerges only over several steps).
This makes it robust to the sign mismatch at deployment and, as a side
effect, prevents it from learning the non-physical through-wall motion
that harms RL training on the MPPI checkpoint. We therefore ship
PPO-WM on the pre\_neg checkpoint and MPPI on the calib196\_160 one;
swapping them degrades PPO-WM to $\sim$60\% sim success (wall-stuck
failures) and is not recommended. See \texttt{hardware/README.md} for the
full sign-convention chain from PLC command to Python state.

\subsection{PPO Training on World Model}

We wrap the world model as a Gymnasium environment~\cite{towers2024gymnasium}
(\texttt{BalancerWorldModel}) and train PPO inside it with 16 parallel
environments (DummyVecEnv) for 3\,M timesteps, using the balanced reward and
no domain randomization. Training takes about 42 minutes on CPU.

\subsection{Results}

\begin{table}[!t]
\centering
\caption{World-model-trained PPO vs.\ physics-trained PPO (50 hardware trials).
  The PPO~(BZ+DR) baseline here is the S-curve development policy; Table~II
  (main paper) reports the matched first-order-lag checkpoint at 96\%
  (Appendix~\ref{appendix:rl_fairness}).}
\label{tab:wm_results}
\begin{tabular}{lcccc}
\toprule
Model & HW SR & HW Mean & HW Med & Sim Settle (S-curve) \\
\midrule
PPO (BZ+DR)           & 98\%   & 5.11\,s & 2.55\,s & 1.67\,s \\
\textbf{PPO (WM)}     & \textbf{100\%} & 6.84\,s & \textbf{3.64\,s} & 1.74\,s \\
\bottomrule
\end{tabular}
\end{table}

The world-model-trained agent achieves \textbf{100\% success rate} on hardware
(50/50 trials), the only RL controller in the deployed 13-controller benchmark
to reach perfect reliability.  The mean settling time is 6.84\,s (std=6.94\,s) but the
median of 3.64\,s is more representative; a few difficult initial
conditions produce outliers up to 26\,s.  In simulation (first-order motor model, 50 trials), it achieves
98\% SR with 1.96\,s mean settling.

\subsection{Data Size Sweep}

To determine how much data the world model actually requires, we trained
the same LSTM architecture on subsets of the random-exploration dataset
ranging from 10k to 500k transitions (the main model used 1M). Open-loop
prediction accuracy (short-horizon, 50 steps) and, for a subset,
10-trial hardware validation under MPPI were evaluated for each model.

\begin{table}[!t]
\centering
\caption{World model prediction error vs.\ training data size.
  Per-channel MAE averaged over the $h=50$ (2.5\,s) rollout.
  The reference analytical model achieves 0.475 total MAE over the same
  rollout.}
\label{tab:wm_data_sweep}
\begin{tabular}{lcccc}
\toprule
\textbf{Transitions} & $x$ (m) & $\dot{x}$ (m/s) & $\theta$ (rad) & $\dot{\theta}$ (rad/s) \\
\midrule
10k  & 0.145 & 0.223 & 0.042 & 0.090 \\
25k  & 0.091 & 0.148 & 0.044 & 0.077 \\
50k  & 0.079 & 0.127 & 0.040 & 0.072 \\
70k  & 0.084 & 0.130 & 0.040 & 0.071 \\
100k & 0.080 & 0.123 & 0.037 & 0.069 \\
150k & 0.075 & 0.117 & 0.039 & 0.070 \\
200k & 0.073 & 0.116 & 0.042 & 0.071 \\
300k & 0.070 & 0.111 & 0.040 & 0.070 \\
400k & 0.074 & 0.118 & 0.038 & 0.068 \\
500k & 0.072 & 0.112 & 0.036 & 0.069 \\
1M   & 0.071 & 0.113 & 0.040 & 0.069 \\
\bottomrule
\end{tabular}
\end{table}

Prediction accuracy plateaus around 50--100k transitions: from 100k
onward the per-channel MAE stays within about $\pm0.01$ of the 1M model
(at table precision), and
scaling to 500k or 1M offers no meaningful improvement.  The 10k model
is the only one noticeably worse (total MAE 0.499); it alone is
comparable to the analytical model's total MAE of 0.475, while every
model from 25k onward clearly outperforms the analytical reference.

To confirm that the prediction plateau translates to controller
performance, we ran 10 hardware trials of MPPI (the same tuning as the
main world-model MPPI) on each trained model.

\begin{table}[!t]
\centering
\caption{Hardware validation of MPPI on world models trained with
different data sizes (10 trials each; the 1M reference row is the first
10 trials of the 50-trial \texttt{MPPI\_v12} hardware run, as no separate
1M sweep run was collected). SR: success rate;
settling time reported as mean and median over successful trials.}
\label{tab:wm_data_sweep_hw}
\begin{tabular}{lccc}
\toprule
\textbf{WM data} & \textbf{HW SR} & \textbf{HW Mean} & \textbf{HW Median} \\
\midrule
50k            & 80\% (8/10)  & 7.03\,s  & 5.11\,s \\
70k            & 100\% (10/10) & 9.65\,s  & 9.90\,s \\
100k           & 70\% (7/10)  & 12.89\,s & 10.43\,s \\
300k           & 100\% (10/10) & 5.30\,s  & 2.85\,s \\
400k           & 100\% (10/10) & 10.18\,s & 6.42\,s \\
500k           & 100\% (10/10) & 8.88\,s  & 8.83\,s \\
1M (main)      & 100\% (10/10) & 1.90\,s  & 1.92\,s \\
\bottomrule
\end{tabular}
\end{table}

The hardware results are noisier than the open-loop errors and only
partly track the plateau: the 70k model reaches 10/10 while the 100k
model reaches 7/10, even though its prediction MAE sits squarely on the
plateau. Each point is a single 10-trial run, and every sweep model runs
under the MPPI tuning optimized for the 1M main model, so a dip like the
100k model's may come from a model-planner mismatch rather than from the
data amount alone. The 1M main model sets the reference (100\% SR,
1.90\,s mean), the lowest settling in the table. The supported finding is
the open-loop one: prediction error, measured on the held-out validation
split, plateaus by 50k to 100k transitions, well below our full
collection. The hardware sweep gives the rough scale of the data
requirement, not a precise threshold.

We also tested a model trained on 70k transitions from a
\emph{balanced demonstration} dataset (collected by a simple heuristic
controller rather than random actions). With the same MPPI tuning as the
main model, it achieves 80\% SR on 10 hardware trials with a mean
settling of 15.3\,s, below the equivalent 70k random-exploration
model (100\% SR, 9.65\,s). At equal data size, this structured-demonstration
model did not match random exploration; richer demonstration data could
narrow the gap, but these single runs do not show that structured collection
lowers the data requirement.


